% ---------------------------------------------------------------------
\subsection{Pipes as bodies in the room (\fn{OrganBodies.cpp})}
\label{sec:math-bodies}
In the room motions (\S\ref{sec:roommotion}) every pipe is a point with a position $\mathbf x$ and a velocity $\mathbf v$
in the Saal.

\paragraph{Integration (\fn{stepBodies}).} The bodies advance in fixed steps of $h = 1/480$\,s, counted in whole samples
and with the remainder carried to the next block, so the paths and the impacts are the same at every host block size
and sample rate (a variable step would make the bounces differ from run to run). Each step is a semi-implicit
(symplectic) Euler step, $\mathbf v \leftarrow \mathbf v + \mathbf a\,h$, then $\mathbf x \leftarrow \mathbf x + \mathbf v\,h$
with the \emph{new} velocity, which, unlike the explicit Euler method, does not gain energy in an orbit or a bounce. A
stalled host never makes the bodies catch up more than 2400 steps (five seconds), and the speed is limited to
$1.5\cdot\SI{40}{m/s}$. The \ui{Geschwindigkeit} $v_0$ of the stop is $10$\,m/s per Hz of its motion rate.

\paragraph{The modes.}
\begin{itemize}[nosep]
  \item \ui{Flug}: no force; the speed returns towards $v_0$, $\mathbf v \leftarrow \mathbf v\bigl(1 + (v_0/|\mathbf v| - 1)\min(1, 2h)\bigr)$,
    and the walls are perfectly elastic.
  \item \ui{Fall}: gravity $a_z = -\SI{9.81}{m/s^2}$, a little air drag ($0.05$ per second horizontally), friction $4$ per
    second while on the floor; a pipe resting longer than $0.8$--\SI{2.3}{s} is thrown up again with $v_z\approx v_0$.
    The coefficient of restitution of floor and walls is $0.3+0.65\,d$ with the depth $d$.
  \item \ui{Schwarm}: Reynolds' three rules of flocking \citep{reynolds1987}: cohesion $0.6(\bar{\mathbf x}-\mathbf x)$
    towards the flock's centre, alignment $1.5(\bar{\mathbf v}-\mathbf v)$ with its mean velocity, and separation
    $0.16(1/r^2 - 1)\,\mathbf r$ from neighbours closer than \SI{1}{m} (evaluated every eighth step per pipe, staggered,
    since it is quadratic in the number of pipes), plus a random push proportional to the depth; the speed is then
    pulled towards $v_0$.
  \item \ui{Strömung}: the pipe relaxes towards the velocity of a slow current
    $\mathbf u(\mathbf x,t)$, made of a few sines of the normalised coordinates and of time (smooth in space and
    time, and divergence-free: $u_x$ depends only on $y,z$, $u_y$ only on $x,z$ and $u_z$ only on $x,y$, so
$\nabla\cdot\mathbf u = 0$ and the current neither gathers nor scatters the pipes), with the rate
    $\min(1,1.5h)$ per step, plus turbulence proportional to the depth.
  \item \ui{Schweben}: a damped spring back to the place where it started,
    $\mathbf a = k(\mathbf x_0 - \mathbf x) - c\,\mathbf v + \text{noise}$, with $k = 2.5(1-d)^2$ and
    $c = 0.35 + 1.6(1-d)$: at full depth nothing calls it back.
\end{itemize}

\paragraph{The player and the hands.} The player attracts or repels the pipes with a softened gravity,
\begin{equation}
  a = \frac{30\,F\,R^2}{r^2 + 0.3R^2},
\end{equation}
(force $F$, range $R$), which is $\approx 30F R^2/r^2$ far away but finite at the player; in the orbiting mode the
tangential velocity is steered towards $v_c = \sqrt{a\,r_h}$, the speed of a circular orbit under that acceleration
($v^2/r = a$), so the pipes circle instead of falling in; the wake mode makes nearby pipes take on the player's velocity
with the weight $6FR^2/(r^2+R^2)$. A carried pipe's velocity relaxes (at the rate 30 per second) towards
$14(\mathbf x_\mathrm{hand} - \mathbf x)$ (at most \SI{25}{m/s}), so it follows the hand closely and keeps the hand's speed
when let go: it is thrown. Arranging moves a pipe by the spring $16\,\Delta\mathbf x - 8\mathbf v$ (critically damped:
$c^2 = 4k$), calling home by $9\,\Delta\mathbf x - 6\mathbf v$ (also critical, $36 = 4\cdot9$), so both arrive without
overshooting. A whirlwind sets the tangential speed $6(1 - r/7)$ within \SI{7}{m}, with a slight inward pull and a lift
$\propto(1-r/7)^2$.

\paragraph{Collisions.} A pipe stands with its whole length between floor and ceiling ($z\in[\ell/2, H-\ell/2]$) and
keeps \SI{0.25}{m} from the walls. At a wall with the inward normal $\mathbf n$ the normal velocity $v_n = \mathbf v\cdot\mathbf n<0$ is
reflected with the restitution $e$, $\mathbf v \leftarrow \mathbf v - (1+e)\,v_n\mathbf n$, and the overshoot of the position is
mirrored back by the factor $e$ as well. The speed of the hit $|v_n|$ and the absorption of the surface are passed to
the knock.

% ---------------------------------------------------------------------
\subsection{The knock (\fn{startImpact}, \fn{renderImpacts})}
\label{sec:math-impact}
An impact is synthesised by modal synthesis: a short contact force excites a few resonators that stand for the struck
pipe's modes, plus a little noise and the thump of the wall.

\paragraph{The contact.} With the speed of the hit $v$ and the absorption $\alpha$ of the surface, the level is
$g = 0.045\,L\,(v/3)^{1.25}(1-0.65\alpha)$ (louder with speed somewhat faster than proportionally; a soft wall takes much
of it) and the contact lasts
\begin{equation}
  T_c = \Bigl(0.0003 + \frac{0.003}{1+v}\Bigr)(1+3\alpha)\ \mathrm{s}
\end{equation}
(three times longer for the wooden Gedackt): a fast hit on a hard surface is short. The force is a raised-cosine pulse
of length $T_c$ and area $g$. Its spectrum is flat up to about $1/T_c$ and falls off above, so a short contact excites high
modes and a long one does not; the code scales each mode of frequency $f$ by $1/(1+(fT_c)^2)$, which is this low-pass
character in its simplest form. Noise low-passed at $f_c = (700 + 11000\min(1,v/6))(1-0.7\alpha)$\,Hz is added during the
contact (the clatter of a fast hit).

\paragraph{The modes of the tube.} A pipe hit from the side rings like a free bar in bending. The Euler--Bernoulli beam
with free ends has its modes at the ratios $1 : 2.756 : 5.404 : 8.933$ (the squares of the roots of
$\cos\beta\cosh\beta = 1$, relative to the first). The fundamental is taken as $f_1 = 300\,\ell^{-1.5}$\,Hz for a
pipe of length $\ell$ (0.7 times that for wood), so long pipes ring low. Each mode is a two-pole resonator,
\begin{equation}
  y_n = 2R\cos\omega\,y_{n-1} - R^2y_{n-2} + g_m\sin\omega\cdot x_n,\qquad R = e^{-6.91/(T_{60}f_s)},
\end{equation}
whose impulse response is $R^n\sin((n+1)\omega)/\sin\omega$; the factor $\sin\omega$ makes its amplitude independent of the
frequency. The decay times fall with the mode number ($T_{60}\cdot 1, 0.7, 0.5, 0.35$; wood rings only $0.12$ as long, brass
$1.2$ times), and a last resonator at the pipe's own note sounds faintly ($0.25$), since a struck pipe also speaks a
little. The Gedackt gets a dry resonance of its hollow box, $700\sqrt{0.05/r}$\,Hz with $T_{60} = \SI{70}{ms}$. Unless a
pipe hit a pipe, the wall itself adds a thump: noise through a one-pole at 110\,Hz (floor) or \SI{180}{Hz} (wall) with
$\tau = \SI{30}{ms}$. The knock is placed where it happened, seen from the listener, and sent to the room.

% ---------------------------------------------------------------------
\subsection{Beats and difference tones (\fn{updateBeats})}
\label{sec:math-beats}
For every pair of sounding pipes $f_1\le f_2$ the engine looks for the lowest pair of partials that nearly coincide:
for $m = 1\dots6$ it takes $n = \mathrm{round}(mf_2/f_1)\le8$ and the beat $b = |nf_1 - mf_2|$, keeping the pair with the
smallest product $nm$ whose beat is below \SI{35}{Hz} (the simplest ratio that beats audibly; the product $nm$ is a
common measure of how simple a ratio is). Every pipe takes the fastest beat it is part of, smoothed with
$\tau = \SI{250}{ms}$, and the source \ui{Schwebung} is $\log_2(1+b)/\log_2 33$ (32\,Hz $=1$). In the motion
\ui{Schwebung} a pipe turns by $360^\circ\cdot\rho\,b\,\Delta t$ per block: its beat becomes its speed around the
listener, so a pure chord stands still and an impure one spins.

Difference tones (\ui{Geister}) are $f_d = f_2 - f_1$ for pairs with $16\,\mathrm{Hz}\le f_d\le 1.02f_1$, below the lower
pipe, as the resultant bass of organ builders; their weight is $\sqrt{g_1g_2}$ times a blend between 1 and the
consonance $\min(1, 6/nm)$, and the strongest few are played.

% ---------------------------------------------------------------------
\subsection{Wandering echoes (\fn{renderEchoes})}
\label{sec:math-echoes}
The dry organ ($W$, or the mid of the stereo pair) is written to a ring buffer. Repeat $r = 1\dots N$ reads it at the
delay $r\,D$ (with $D$ in seconds or in note values of the tempo; a changed $D$ glides with $\tau = \SI{120}{ms}$, read with
linear interpolation, which makes the glide a tape-like pitch bend instead of a click), with the gain
$G\,\phi^{r-1}$ (level $G$, \ui{Abnahme} $\phi$) and a one-pole low-pass at
$\max\bigl(250,\ 18000\,(1-0.85\delta)^r\bigr)$\,Hz, darker each time. It is encoded from the direction
$(r\,\Delta\varphi,\ r\,\Delta\vartheta)$, turning by \ui{Drehung} and tilting by \ui{Neigung} per repeat, and also sent to
the room from that direction with the circular-harmonic weights of \S\ref{sec:math-ir}. All repeats read the same
buffer: there is no feedback loop, so the echoes cannot build up, and each one can come from its own place.
